\section*{Hoofdstuk \ref{ch:b1:classifying-biodiversity} --- De levende wereld indelen}

\begin{solution}{exo:b1:classifying-biodiversity:1}
\emph{Felis catus}: \omterm{def:b1:classifying-biodiversity:nomenclature}{geslacht} \emph{Felis}, soortaanduiding \emph{catus};
familie Felidae, orde Carnivora, klasse Mammalia, stam Chordata.
\end{solution}

\begin{solution}{exo:b1:classifying-biodiversity:2}
\omterm{def:b1:classifying-biodiversity:characters}{Homologie}: gelijkenis door afstamming van een gemeenschappelijke voorouder
(de botten van de voorpoot van een vleermuis en van een mens). \omterm{def:b1:classifying-biodiversity:characters}{Homoplasie}:
gelijkenis door convergentie of terugkeer (de vleugels van een vleermuis en
van een vogel). Alleen gedeelde afgeleide homologieën zijn bewijs voor een
\omterm{def:b1:classifying-biodiversity:systematics}{clade}.
\end{solution}

\begin{solution}{exo:b1:classifying-biodiversity:3}
``Vissen'': parafyletisch (sluit de viervoeters uit die van vissen
afstammen). Vogels: monofyletisch (alle nakomelingen van één bevederde
voorouder). ``Warmbloedige dieren'': polyfyletisch (vogels en zoogdieren
kregen de endothermie apart).
\end{solution}

\begin{solution}{exo:b1:classifying-biodiversity:4}
Bacteria, Archaea, Eukarya. Archaea hebben etergebonden membraanlipiden en
geen \omterm{def:b1:carbohydrates:polysaccharide}{peptidoglycaan}; hun \omterm{def:b1:gene-expression:transcription}{RNA-polymerase} en ribosomale \omterm{def:b1:proteins:peptide}{eiwitten} lijken op die
van de eukaryoten.
\end{solution}

\begin{solution}{exo:b1:classifying-biodiversity:5}
Bacteriën planten zich niet onderling voort --- zij delen zich, en wisselen
\omterm{def:b1:genomes:gene}{genen} uit over wat soortsgrenzen zouden zijn --- en fossielen kunnen niet op
onderlinge voortplanting worden getoetst. Bacteriën worden afgebakend met de
gelijkenis van hun volgorden (een drempel op het ribosomale \omterm{def:b1:nucleic-acids:chain}{RNA} of op
genomen); fossielen met hun bouw.
\end{solution}

\begin{solution}{exo:b1:classifying-biodiversity:6}
Boom (X,Y),Z: de kenmerken 1, 3 en 4 elk één verandering, kenmerk 2 twee
veranderingen: lengte 5. Boom (Y,Z),X: kenmerk 2 één verandering, de
kenmerken 1, 3 en 4 elk twee: lengte 7. (X,Y) verdient de voorkeur; kenmerk 2
is de \omterm{def:b1:classifying-biodiversity:characters}{homoplasie}.
\end{solution}

\begin{solution}{exo:b1:classifying-biodiversity:7}
De \omterm{def:b1:classifying-biodiversity:characters}{buitengroep} legt de polariteit van elk kenmerk vast: haar toestand wordt
als oorspronkelijk genomen, zodat alleen de afgeleide toestanden als bewijs
tellen. Een \omterm{def:b1:classifying-biodiversity:characters}{buitengroep} die in werkelijkheid binnen de groep ligt laat
sommige afgeleide toestanden oorspronkelijk lijken en wortelt de boom op de
verkeerde plaats, waardoor \omterm{def:b1:classifying-biodiversity:systematics}{claden} parafyletische groepen worden.
\end{solution}

\begin{solution}{exo:b1:classifying-biodiversity:8}
Ribosomale en eiwitvolgorden plaatsen dieren en schimmels als zusterlijnen,
beide binnen de opisthokonten, waarvan de zwemmende \omterm{def:b1:cell-unit-of-life:cell}{cellen} (zaadcellen,
zoösporen van chytriden) één achterwaartse zweepstaart hebben; planten liggen
in een aparte lijn die door de primaire chloroplast wordt bepaald.
\end{solution}

\begin{solution}{exo:b1:classifying-biodiversity:9}
Een rijk binnen de bacteriën zou inhouden dat de methanogenen dichter bij de
overige bacteriën staan dan bij de eukaryoten; de volgorden toonden drie
lijnen die even ver van elkaar af stonden, zodat de prokaryoten geen één
groep waren en de diepste scheiding van het leven drievoudig was.
\end{solution}

\begin{solution}{exo:b1:classifying-biodiversity:10}
Gastheerspecifieke kevers per boomsoort: $0.2\times 1200 = 240$; insecten van
de kruin die eigen zijn aan de boom: $240/0.4 = 600$; hele boom:
$600/(2/3) = 900$; maal \num{50000} bomen: \num{45} miljoen. Aannamen: de
benevelde boom is doorsnee, de specificiteit is echt en geen toevalsfout,
kevers zijn op elke boom \qty{40}{\%} van de insecten, en \omterm{def:b1:classifying-biodiversity:species}{soorten} die aan
één boom eigen zijn worden niet dubbel geteld --- elk daarvan is op een
factor twee of meer onzeker, en latere schattingen zijn naar \num{5} miljoen
gezakt.
\end{solution}

\begin{solution}{exo:b1:classifying-biodiversity:11}
Het \omterm{def:b1:classifying-biodiversity:nomenclature}{typemateriaal} is het voorwerp waarnaar de naam verwijst: de levende
\omterm{def:b1:populations:population}{populaties} zijn dezelfde \omterm{def:b1:classifying-biodiversity:species}{soort} als zij met het type soortgelijk worden
geoordeeld, en een nieuwe \omterm{def:b1:classifying-biodiversity:species}{soort} zo niet; in het tweede geval blijft de oude
naam bij het type en hebben de nieuwe \omterm{def:b1:populations:population}{populaties} een nieuwe naam nodig. De
voorrang waarborgt dat wanneer blijkt dat twee namen naar één \omterm{def:b1:classifying-biodiversity:species}{soort}
verwijzen, de oudste wordt gehouden, zodat namen zich niet vermenigvuldigen
of met elke nieuwe opvatting verschuiven.
\end{solution}

\begin{solution}{exo:b1:classifying-biodiversity:12}
\omterm{def:b1:classifying-biodiversity:parsimony}{Spaarzaamheid} veronderstelt alleen dat convergenties en terugkeren zeldzamer
zijn dan gedeelde afstamming, zodat een boom er niet meer van mag vragen dan
de gegevens afdwingen. Zij faalt wanneer convergentie veel voorkomt
(gelijksoortige \omterm{def:b1:ecosystem-organization:niche}{leefgebieden}), wanneer kenmerken samenhangen (één aanpassing
die als vele wordt geteld), en wanneer snel evoluerende lijnen toevallige
gelijkenissen opstapelen en naar elkaar toe worden getrokken (``lange takken
trekken aan''). Moleculaire volgorden leveren vele onafhankelijke kenmerken
en substitutiemodellen om dezelfde boom te toetsen; fossielen leveren de
volgorde waarin de afgeleide toestanden verschijnen en dateringen voor de
splitsingen; overeenstemming tussen onafhankelijke gegevens is de toets.
\end{solution}

\begin{solution}{pb:b1:classifying-biodiversity:1}
\textbf{1.} Kenmerk 1, kaken: gedeeld door alle \omterm{def:b1:classifying-biodiversity:nomenclature}{taxa} behalve de prik, zodat
het de binnengroep bepaalt en de prik als \omterm{def:b1:classifying-biodiversity:characters}{buitengroep} bevestigt.
\textbf{2.} 1 (7 \omterm{def:b1:classifying-biodiversity:nomenclature}{taxa}: gewervelden met kaken), 2 (6: benige gewervelden), 3
(5: longvis en viervoeters), 4 (4: viervoeters), 5 (3: amnioten), 8, 9
en 10 (elk 2: hagedis+duif, muis+duif, hagedis+duif), 6 en 7
(elk 1).
\textbf{3.} Haren (muis) en veren (duif): unieke toestanden,
autapomorfieën, die het \omterm{def:b1:classifying-biodiversity:nomenclature}{taxon} kenschetsen maar niets over zijn verwanten
zeggen.
\textbf{4.} Kenmerk 9 (endothermie: muis+duif) botst met 8 en 10
(hagedis+duif); de botsing speelt tussen de drie amnioten.
\textbf{5.} De afwezigheid van \omterm{def:b1:gas-exchange:lung}{longen} is oorspronkelijk (de prik, de haai en
de forel missen ze), de aanwezigheid afgeleid; de toestand van de
\omterm{def:b1:classifying-biodiversity:characters}{buitengroep} wordt als oorspronkelijk genomen.
\textbf{6.} Drie: elk van de drie kan de buitenstaander zijn ---
(hagedis,duif) met de muis erbuiten, (muis,duif) met de hagedis erbuiten,
(hagedis,muis) met de duif erbuiten. Alleen de kenmerken 8, 9 en 10
verschillen tussen de drie, dus alleen zij doen ertoe voor de keuze.
\textbf{7.} (prik,(haai,(forel,(longvis,(kikker,(hagedis, muis,
duif)))))), met elke knoop door één kenmerk gemarkeerd.
\textbf{8.} Boom A: (muis,(hagedis,duif)); boom B: (hagedis,(muis,
duif)).
\textbf{9.} Boom A: de kenmerken 1--8 en 10 elk één verandering (9), kenmerk
9 twee veranderingen: lengte \num{11}.
\textbf{10.} Boom B: de kenmerken 1--7 en 9 elk één (8), de kenmerken 8 en
10 elk twee: lengte \num{12}.
\textbf{11.} Boom A, met één stap.
\textbf{12.} De endothermie: zij is twee keer ontstaan, eenmaal in de lijn van
de zoogdieren en eenmaal in die van de vogels --- een convergentie, die
strookt met hun verschillende mechanismen (vacht en zweet tegenover veren en
luchtzakken).
\textbf{13.} Zonder 8 en 10: boom A lengte 9 (endothermie twee keer), boom B
lengte 8 (alles in overeenstemming): B zou de voorkeur krijgen. De conclusie
hangt af van de gekozen kenmerken; met weinig kenmerken kan één \omterm{def:b1:classifying-biodiversity:characters}{homoplasie} de
uitkomst doen kantelen, en daarom zijn er vele kenmerken nodig, en ook
moleculaire.
\textbf{14.} Nee: hun gemeenschappelijke voorouder is ook de voorouder van de
viervoeters, dus alle vier de viervoeters zouden erbij moeten.
\textbf{15.} Ja: het amnionei.
\textbf{16.} Nee: de duif stamt van binnen de reptielen af (zij is hier de
zuster van de hagedis, en in vollediger bomen van de krokodillen), zodat
reptielen zonder vogels parafyletisch zijn.
\textbf{17.} Nee: hun naaste gemeenschappelijke voorouder is de voorouder van
alle amnioten, die niet \omterm{prop:b1:organism-environment:thermo}{endotherm} was; polyfyletisch, samengebracht op een
convergent kenmerk.
\textbf{18.} Dichter bij de kikker: hij deelt de \omterm{def:b1:gas-exchange:lung}{longen} (kenmerk 3) met de
viervoeters, die de forel mist; ``vissen'' zonder de viervoeters is
parafyletisch.
\textbf{19.} Klasse Reptilia, klasse Mammalia, klasse Aves: drie klassen van
gelijke rang, hoewel de splitsing tussen hagedis en duif veel recenter is dan
die van elk van beide met de zoogdieren; rangen zijn afspraken, en een
gelijke rang betekent geen gelijke ouderdom of gelijke eigenheid.
\textbf{20.} Zoogdieren: \num{20} nieuwe per jaar op \num{6500}
(\qty{0.3}{\%}); vissen: \num{300} op \num{35000} (\qty{0.9}{\%}) --- de
inventaris van de zoogdieren is het dichtst bij voltooiing, die van de vissen
groeit nog snel.
\textbf{21.} Negen tiende van de volgorden hoort bij \omterm{def:b1:organism-environment:organism}{organismen} die nooit zijn
gekweekt of benoemd: de beschreven prokaryoten zijn een klein deel van die in
zelfs één liter water, en de inventaris is nauwelijks begonnen.
\textbf{22.} Achtergrond: \num{9} per jaar. Duizendvoudig: \num{9000} per
jaar --- de meeste ervan \omterm{def:b1:classifying-biodiversity:species}{soorten} die nooit zijn beschreven.
\textbf{23.} Haar kenmerken, haar volgorde en haar plaats op de boom gaan met
haar verloren; een tak van de levensboom wordt uitgewist voordat hij is
getekend, en elke \omterm{def:b1:classifying-biodiversity:systematics}{clade} die alleen zij zou hebben bepaald gaat mee.
\textbf{24.} Families zijn groot, opvallend en vrijwel alle beschreven
(nieuwe families zijn nu zeldzaam), zodat hun telling bijna volledig is; de
regelmatige verhouding van \omterm{def:b1:classifying-biodiversity:species}{soorten} tot \omterm{def:b1:classifying-biodiversity:nomenclature}{geslachten} tot families over de goed
bekende groepen laat de familietelling dan doortrekken naar \omterm{def:b1:classifying-biodiversity:species}{soorten} waar de
soortentelling niet volledig is.
\textbf{25.} (prik,(haai,(forel,(longvis,(kikker,(muis,(hagedis,
duif))))))), lengte \num{11}.
\end{solution}
